\documentclass{../shipyard-slide}

\usepackage[british]{babel}
\usepackage[british]{datetime2}
\usepackage{amsmath}
\usepackage{amssymb}
\usepackage{ulem}
\usepackage{qrcode}
\usepackage{twemojis}
\usepackage{wasysym}

\usepackage{tikz,calc}
\usetikzlibrary{shapes, shapes, arrows, chains, fit, quotes}

\tikzstyle{trienode} = [draw=text-color, rounded corners]
\tikzstyle{bucket} = [draw=text-color, rectangle]

\settheme{shipyard-2025}

%Information to be included in the title page:
\title{Introduction to IPFS}
\subtitle{Web Verifiability through Content Addressing}
\author{Guillaume Michel}
\avatar{../resources/avatar.jpg}
\handle{@guissou}
\institute{Interplanetary Shipyard}
\event{Privacy x Verifiability Conference - EPFL}
\date{\DTMdate{2025-03-07}}

\begin{document}

\frame{\titlepage}

\logo{
	\vspace{7.2cm}
	\includegraphics[scale=.16]{../shipyard-resources/shipyard-symbol-white-trans-bg.png}
	\hspace{4em}
}

\begin{frame}
  \frametitle{Why verifiability matters}
  \begin{itemize}
    \itemc Trust \& Transparency
    \itemc Security \& Accountability
    \itemc DApps Resilience
  \end{itemize}

  \begin{columns}
  \begin{column}{.63\textwidth}
  \begin{itemize}
    \item[\greencube] Recent example: ByBit hack
  \end{itemize}
  \end{column}
  \begin{column}{.30\textwidth}
    \begin{center}
    \qrcode{https://blog.ipfs.tech/2025-could-ipfs-prevent-bybit-hack/}\\
    \medskip
    {\footnotesize \textit{Could IPFS have prevented the Bybit hack?}}
  \end{center}
  \end{column}
  \end{columns}

%       Set the Stage:
%         Traditional web addresses data by location (e.g., URLs).
%         This approach makes content integrity and resilience more challenging (links can break, data can be changed without notice).
%         Growing concerns about misinformation, censorship, and reliability of web content.
%     Introduce the Concept of Verifiability:
%         The ability to cryptographically verify data.
%         Ensures that what you’re retrieving is exactly what was intended by the publisher.
%
% Tip: Use real-world examples (e.g., content takedowns, link rot, or ephemeral data) to highlight the importance of verifiability.
\end{frame}

\begin{frame}
  \frametitle{Agenda}
  \begin{enumerate}
    \item What is IPFS?
    \item Content Addressing
    \item Content Routing
    \item IPFS on the Web
    \item Common Misconceptions
  \end{enumerate}
\end{frame}

% \begin{frame}
% \frametitle{What is IPFS?}
% \begin{columns}
% \begin{column}{0.59\textwidth}
%   \begin{itemize}
%     \itemc IPFS = \textbf{I}nter\textbf{P}lanetary \textbf{F}ile\textbf{S}ystem
%     \itemc Peer-to-peer hypermedia protocol for content addressing
%     \itemc Aims to make the web faster, safer, and more open
%     \itemc Does IPFS have a blockchain? No.
%     \begin{itemize}
%       \item \url{https://doesipfshaveablockchain.com}
%     \end{itemize}
%   \end{itemize}
% \end{column}
% \begin{column}{0.35\textwidth}
%   \textbf{Key Properties}
%   \begin{enumerate}
%     \item Decentralized
%     \item Content Addressing
%     \item Merkle DAGs
%     \item Resilience
%     \item Censorship Resistance
%   \end{enumerate}
% \end{column}
% \end{columns}
%
% \end{frame}

\begin{frame}
\frametitle{What is IPFS?}
\begin{columns}
  \begin{column}{.7\linewidth}
  \begin{itemize}
    \itemc IPFS = \textbf{I}nter\textbf{P}lanetary \textbf{F}ile\textbf{S}ystem
    \itemc Peer-to-peer hypermedia protocol for content addressing
    \itemc Aims to make the web faster, safer, and more open
    \itemc Does IPFS have a blockchain? No.
    \begin{itemize}
      \item \url{https://doesipfshaveablockchain.com}
    \end{itemize}
  \end{itemize}
\end{column}
\begin{column}{.25\linewidth}
  \includegraphics[width=\linewidth]{resources/ipfs-logo.png}
\end{column}
\end{columns}
\end{frame}

% \begin{frame}
%   \frametitle{IPFS Key Properties}
% \begin{enumerate}
%   \item Decentralized
%   \item Content Addressing
%   \item Merkle DAGs
%   \item Resilience
%   \item Censorship Resistance
% \end{enumerate}
%
% \end{frame}

\begin{frame}
  \frametitle{Location Addressing vs Content Addressing}
  
\begin{center}
  \renewcommand{\arraystretch}{1.5}
  {\small
  \begin{tabular}{|p{.45\linewidth}|p{.45\linewidth}|}
  \hline
  \large\textbf{Location Addressing} & \large\textbf{Content Addressing} \\ \hline
  \textit{Where is the data} & \textit{What is the data?} \\ \hline
  Uniform Resource Locator (URL) & Content Identifier (CID) \\
  \texttt{https://example.com/foo/bar.jpg} & \texttt{bafybeihfg3d7rdltd43u3tfvn...uc7y} \\ \hline
  Server down? The data is unreachable & Anyone can serve the same data \\ \hline
  Limited verifiability — trust the server to provide correct data & Full verifiability — if content changes, CID changes \\ \hline
  Censor-prone — block IP or domain name & Hard to censor — all providers must be blocked \\ \hline
  \end{tabular}
  }
\end{center}
\end{frame}

\begin{frame}
  \frametitle{Content Identifiers (CIDs)}
  \begin{itemize}
    \itemc Immutable, self certifying name
    \itemc Made of hash and meta data
  \end{itemize}
  \bigskip
  \includegraphics[width=\linewidth]{resources/cid-composition.png}
  {\small source: \url{cid.ipfs.tech}}
\end{frame}


\begin{frame}
  \frametitle{Merkle DAG}
  \begin{center}
  \includegraphics[width=.7\linewidth]{resources/file.png}
  \end{center}
\end{frame}
\begin{frame}
  \frametitle{Merkle DAG}
  \begin{center}
  \includegraphics[width=.7\linewidth]{resources/chunks.png}
  \end{center}
\end{frame}
\begin{frame}
  \frametitle{Merkle DAG}
  \begin{center}
  \includegraphics[width=.7\linewidth]{resources/dag.png}
  \end{center}
\end{frame}
\begin{frame}
  \frametitle{Merkle DAG}
  \begin{center}
  \includegraphics[width=.7\linewidth]{resources/cid.png}
  \end{center}
\end{frame}

\begin{frame}
  \frametitle{libp2p}
  \begin{itemize}
    \itemc Open source peer-to-peer networking library
    \itemc Peer ID: unique peer identifier derived from public key (\texttt{12D3KooWN1Mr...})
    \itemc Transports: \texttt{tcp}, \texttt{quic}, \texttt{wss}, \texttt{webrtc}, \texttt{webtransport}
    \itemc Encrypted by default
    \itemc Features: Holepunching, Relays, Kademlia DHT, Gossipsub, etc.
    \itemc Networking layer of IPFS, Ethereum, Polkadot, Filecoin, etc.
  \end{itemize}
  \begin{tikzpicture}[remember picture, overlay]
    \node[anchor=south east, xshift=-0.7cm, yshift=0.5cm] at (current page.south east)
    {\includegraphics[width=.15\linewidth]{resources/libp2p.png}};
  \end{tikzpicture}
\end{frame}


\begin{frame}
  \frametitle{Content Routing}
  \begin{center}
    \includegraphics[width=.7\textwidth]{resources/routing-add.png}
  \end{center}
\end{frame}
\begin{frame}
  \frametitle{Content Routing}
  \begin{center}
    \includegraphics[width=.7\textwidth]{resources/routing-provide.png}
  \end{center}
\end{frame}
\begin{frame}
  \frametitle{Content Routing}
  \begin{center}
    \includegraphics[width=.7\textwidth]{resources/routing-lookup-bob.png}
  \end{center}
\end{frame}
\begin{frame}
  \frametitle{Content Routing}
  \begin{center}
    \includegraphics[width=.7\textwidth]{resources/routing-peer-lookup.png}
  \end{center}
\end{frame}
\begin{frame}
  \frametitle{Content Routing}
  \begin{center}
    \includegraphics[width=.7\textwidth]{resources/routing-transfer-firewall.png}
  \end{center}
\end{frame}
\begin{frame}
  \frametitle{Content Routing}
  \begin{center}
    \includegraphics[width=.7\textwidth]{resources/routing-verification.png}
  \end{center}
\end{frame}
\begin{frame}
  \frametitle{Content Routing}
  \begin{center}
    \includegraphics[width=.7\textwidth]{resources/routing-provide-bob.png}
  \end{center}
\end{frame}
\begin{frame}
  \frametitle{Content Routing}
  \begin{center}
    \includegraphics[width=.7\textwidth]{resources/routing-lookup-charlie.png}
  \end{center}
\end{frame}
\begin{frame}
  \frametitle{Content Routing}
  \begin{center}
    \includegraphics[width=.7\textwidth]{resources/routing-transfer-charlie.png}
  \end{center}
\end{frame}
\begin{frame}
  \frametitle{Content Routing}
  \begin{center}
    \includegraphics[width=.7\textwidth]{resources/routing-no-providers.png}
  \end{center}
\end{frame}
\begin{frame}
  \frametitle{Content Routing}
  \begin{center}
    \includegraphics[width=.7\textwidth]{resources/routing-ipni.png}
  \end{center}
\end{frame}


\begin{frame}
  \frametitle{Bridging the Web: IPFS HTTP Gateways}
  \begin{itemize}
    \itemc Browsers have limited connectivity:
    \begin{itemize}
      \item[\greencube] \texttt{wss}
      \item[\greencube] \texttt{webrtc}
      \item[\greencube] \texttt{webtransport} (WIP)
    \end{itemize}
    \itemc HTTP can be used for data retrieval!
  \end{itemize}
\end{frame}
\begin{frame}
  \frametitle{Bridging the Web: IPFS HTTP Gateways}
  \begin{center}
    \includegraphics[width=.7\textwidth]{resources/routing-gateway.png}
  \end{center}
\end{frame}
\begin{frame}
  \frametitle{Bridging the Web: IPFS HTTP Gateways}
  \begin{itemize}
    \itemc HTTP Gateways (\href{https://ipfs.io}{ipfs.io}, \href{https://dweb.link}{dweb.link})
    \begin{itemize}
      \item[\greencube] \texttt{GET /ipfs/\{cid\}}
      \item[\greencube] Gateway must be trusted
    \end{itemize}
    \itemc Trustless gateways (\url{inbrowser.dev})
    \begin{itemize}
      \item[\greencube] Data verified in Service Worker
      \item[\greencube] Uses Helia verified fetch
    \end{itemize}
    \itemc Next Steps
    \begin{itemize}
      \item[\greencube] Delegated Routing
      \item[\greencube] Fetch from source
    \end{itemize}
  \end{itemize}
\end{frame}

\begin{frame}
  \frametitle{Common Misconceptions}
  \begin{columns}
    \begin{column}{.38\textwidth}
  \begin{enumerate}
    \item Uploading files to IPFS
    \item Permanent data storage
    \item Data replication
    \item IPFS fully replaces HTTP
    \item Privacy
  \end{enumerate}
\end{column}
\begin{column}{.5\textwidth}
  \includegraphics[width=\textwidth]{resources/ipfs-upload.png}
\end{column}
\end{columns}
\end{frame}

\begin{frame}
  \frametitle{Shipyard}
  \begin{columns}
    \begin{column}{.8\textwidth}
  \begin{itemize}
    \itemc Engineering collective stewarding IPFS \& libp2p
    \begin{itemize}
      \item[\greencube] We \twemoji{red heart} Open Source
      \item[\greencube] Graduated from Protocol Labs in 2024
    \end{itemize}
    \itemc Championing IPFS \& libp2p specs
    \itemc Developping \& maintaining \href{https://github.com/ipfs/kubo}{kubo}, \href{https://github.com/ipfs/helia}{helia}, \href{https://github.com/libp2p/go-libp2p}{go-libp2p}, \href{https://github.com/libp2p/js-libp2p}{js-libp2p} and associated tooling
    \itemc Operating IPFS Gateways (\href{https://ipfs.io}{ipfs.io}) and DHT Bootstrappers
    \itemc Open for collaborations and supporting your IPFS use cases
    \begin{itemize}
      \item[\greencube] \href{mailto:contact@ipshipyard.com}{contact@ipshipyard.com}
    \end{itemize}
  \end{itemize}
  \end{column}
  \begin{column}{.18\textwidth}
    \begin{center}
    \qrcode{https://ipshipyard.com}\\
    \smallskip
    {\tiny \href{https://ipshipyard.com}{https://ipshipyard.com}}\\
    \bigskip
    \includegraphics[width=.7\textwidth]{../shipyard-resources/shipyard-vertical-white-trans-bg.png}
    \end{center}
  \end{column}
  \end{columns}
  \end{frame}

\begin{frame}
\frametitle{Q\&A}

\begin{columns}[onlytextwidth]
	\begin{column}{0.69\textwidth}
	\begin{columns}[onlytextwidth]
	\begin{column}{.2\textwidth}
		\tikz\node [circle, minimum width = \linewidth,
			path picture = {
      			\node [] at (path picture bounding box.center) {
        				\includegraphics[width=\linewidth]{../resources/avatar.jpg}
        			};
    		}] {};
	\end{column}
	\begin{column}{.78\textwidth}
		{\Large Guillaume Michel (@guissou)}
	\end{column}
	\end{columns}
	\vspace{1cm}
	\begin{itemize}
    \itemc Github: \href{https://github.com/guillaumemichel}{\textbf{@guillaumemichel}}
    \itemc Email: \href{mailto:guillaume@ipshipyard.com}{guillaume@ipshipyard.com}
    \itemc Website: \url{https://ipshipyard.com}
	\end{itemize}

	\end{column}
	\begin{column}{0.29\textwidth}
		\begin{center}
			\qrcode[height=3cm]{https://ipfs.io/ipfs/bafybeid4pnwg2djmlq32qvgnwczswff5aouzrnm65d4bprjxbx22wjxmw4}\\
			\smallskip
			{\small Slides available on IPFS}
		\end{center}
	\end{column}
\end{columns}

\end{frame}
\end{document}
